\documentclass[10pt,a4paper]{amsart}
\usepackage{amsmath,amssymb,amsthm,graphicx}
% Retained minimal preamble for the Stacks Project chapter more-algebra.tex (More on Algebra).
% The chapter file pinned-original-source.tex carries only \input{preamble} on its first line, so no
% preamble text lives inside the chapter itself.  The Stacks Project uses one single shared preamble for
% every chapter; the snapshot stages it identically as sources/stacks-chow/preamble.tex and
% sources/stacks-etale-cohomology/preamble.tex.  The block below is copied mechanically, line for line,
% from sources/stacks-chow/preamble.tex lines 14-18 (the Xy-pic setup) and 154-217 (the native theorem
% declarations together with the mathematical macros).  Nothing was added, removed or redefined; the
% wrapper class and the amsmath/amssymb/amsthm/graphicx base packages are supplied by the assembler.
% Chapter: more-algebra.tex -- More on Algebra.
% For commutative diagrams we use Xy-pic
\usepackage[all]{xy}

% We use 2cell for 2-commutative diagrams.
\xyoption{2cell}

% Theorem environments.
%
\theoremstyle{plain}
\newtheorem{theorem}[subsection]{Theorem}
\newtheorem{proposition}[subsection]{Proposition}
\newtheorem{lemma}[subsection]{Lemma}

\theoremstyle{definition}
\newtheorem{definition}[subsection]{Definition}
\newtheorem{example}[subsection]{Example}
\newtheorem{exercise}[subsection]{Exercise}
\newtheorem{situation}[subsection]{Situation}

\theoremstyle{remark}
\newtheorem{remark}[subsection]{Remark}
\newtheorem{remarks}[subsection]{Remarks}

\numberwithin{equation}{subsection}

% Macros
%
\def\lim{\mathop{\mathrm{lim}}\nolimits}
\def\colim{\mathop{\mathrm{colim}}\nolimits}
\def\Spec{\mathop{\mathrm{Spec}}}
\def\Hom{\mathop{\mathrm{Hom}}\nolimits}
\def\Ext{\mathop{\mathrm{Ext}}\nolimits}
\def\SheafHom{\mathop{\mathcal{H}\!\mathit{om}}\nolimits}
\def\SheafExt{\mathop{\mathcal{E}\!\mathit{xt}}\nolimits}
\def\Sch{\mathit{Sch}}
\def\Mor{\mathop{\mathrm{Mor}}\nolimits}
\def\Ob{\mathop{\mathrm{Ob}}\nolimits}
\def\Sh{\mathop{\mathit{Sh}}\nolimits}
\def\NL{\mathop{N\!L}\nolimits}
\def\CH{\mathop{\mathrm{CH}}\nolimits}
\def\proetale{{pro\text{-}\acute{e}tale}}
\def\etale{{\acute{e}tale}}
\def\QCoh{\mathit{QCoh}}
\def\Ker{\mathop{\mathrm{Ker}}}
\def\Im{\mathop{\mathrm{Im}}}
\def\Coker{\mathop{\mathrm{Coker}}}
\def\Coim{\mathop{\mathrm{Coim}}}

% Boxtimes
%
\DeclareMathSymbol{\boxtimes}{\mathbin}{AMSa}{"02}

%
% Macros for moduli stacks/spaces
%
\def\QCohstack{\mathcal{QC}\!\mathit{oh}}
\def\Cohstack{\mathcal{C}\!\mathit{oh}}
\def\Spacesstack{\mathcal{S}\!\mathit{paces}}
\def\Quotfunctor{\mathrm{Quot}}
\def\Hilbfunctor{\mathrm{Hilb}}
\def\Curvesstack{\mathcal{C}\!\mathit{urves}}
\def\Polarizedstack{\mathcal{P}\!\mathit{olarized}}
\def\Complexesstack{\mathcal{C}\!\mathit{omplexes}}
% \Pic is the operator that assigns to X its picard group, usage \Pic(X)
% \Picardstack_{X/B} denotes the Picard stack of X over B
% \Picardfunctor_{X/B} denotes the Picard functor of X over B
\def\Pic{\mathop{\mathrm{Pic}}\nolimits}
\def\Picardstack{\mathcal{P}\!\mathit{ic}}
\def\Picardfunctor{\mathrm{Pic}}
\def\Deformationcategory{\mathcal{D}\!\mathit{ef}}
\usepackage{cleveref}
\crefname{originaldefinition}{original definition}{original definitions}
\Crefname{originaldefinition}{Original definition}{Original definitions}
\makeatletter
\newlabel{derived-definition-cone}{{014E}{source}{}{sourcepointer-1}{}}
\newlabel{derived-definition-cone@cref}{{[originaldefinition][0][]014E}{[1][0][]source}{}{}{}}
\makeatother
\crefname{originallemma}{original lemma}{original lemmas}
\Crefname{originallemma}{Original lemma}{Original lemmas}
\makeatletter
\newlabel{lemma-cone-pseudo-coherent}{{064R}{source}{}{sourcepointer-2}{}}
\newlabel{lemma-cone-pseudo-coherent@cref}{{[originallemma][0][]064R}{[1][0][]source}{}{}{}}
\makeatother
\crefname{originallemma}{original lemma}{original lemmas}
\Crefname{originallemma}{Original lemma}{Original lemmas}
\makeatletter
\newlabel{lemma-finite-cohomology}{{064S}{source}{}{sourcepointer-3}{}}
\newlabel{lemma-finite-cohomology@cref}{{[originallemma][0][]064S}{[1][0][]source}{}{}{}}
\makeatother
\title{A complex is pseudo-coherent exactly when it is m-pseudo-coherent for all m, and then its finite free resolution can be chosen to vanish above the top cohomology}
\author{The Stacks Project Authors}
\date{stacks-project:more-algebra.tex:local-snapshot-20261012; GNU FDL 1.2 or later}
\begin{document}
\maketitle
\setcounter{subsection}{1}
\noindent\textit{Editorial note.} All delivered mathematical text is the author's own, copied from the pinned Stacks Project chapter file more-algebra.tex; no statement or proof step was rewritten, reordered or removed.\par\medskip
\noindent\textit{Editorial note.} The chapter file pinned-original-source.tex carries only an input of the shared preamble on its first line, so no preamble text lives inside the chapter. The retained file sources/stacks-more-algebra/retained-preamble.tex is a mechanical, line-for-line copy of the project's single shared preamble as staged in the snapshot, namely sources/stacks-chow/preamble.tex lines 14-18 (Xy-pic setup) and 154-217 (theorem declarations and mathematical macros). Nothing was added, removed or redefined, no package was added, and the excerpt serves only as preamble, never as a span coordinate.\par\medskip
\noindent\textit{Editorial note.} The source URL is the source-level Stacks Project URL, as required when more than one card is produced from one source; the per-claim tags are carried by claim\_key and by the external references. The two same-chapter pointers lemma-cone-pseudo-coherent (064R) and lemma-finite-cohomology (064S) are recorded as non-load-bearing navigation stubs only. The definition definition-pseudo-coherent of the same section is shipped neither as a stub nor as a context span because that definition carries a footnote with a bibliography citation and this card includes no bibliography; the terminology is recorded in the allowed standard prerequisites instead.\par\medskip
\noindent\textit{Editorial note.} The proof uses Xy-pic only through the project's own xymatrix macro, supplied by the retained preamble. The same-chapter text immediately preceding the counted span contains that bibliography citation and is deliberately kept outside both counted spans.\par\medskip
\noindent\textit{Allowed standard prerequisites (editorial).}\begin{itemize}
\item The chapter's own definition of an m-pseudo-coherent and of a pseudo-coherent object of D(R), and of the same terms for R-modules. Definition definition-pseudo-coherent at source lines 15873-15893 of the same chapter, referred to by the text at source line 15903. Fixes what the three listed conditions mean: condition (2) asks for m-pseudo-coherence for every m in Z, condition (3) is the reformulation in terms of a bounded above complex of finite projective modules, and the last sentence of the statement refines the resolution that (3) produces.
\item Bounded complexes of finite free or finite projective R-modules and the stupid truncation sigma\_\{>= m\} of a complex. used at proof lines 16064-16076 and 16079-16084. Supplies the objects the three conditions speak about and the truncation that turns a bounded above finite free resolution into a witness of m-pseudo-coherence for a single m.
\item The cone of a morphism of complexes of R-modules together with the long exact cohomology sequence of the resulting distinguished triangle. used at proof lines 16102-16112, where the cone is the one of Derived Categories, Definition derived-definition-cone. Converts the two prescribed cohomological properties of the partially built map alpha into the vanishing of H\textasciicircum{}i(C\textasciicircum{}bullet) for i >= n, which is what lets the induction step be restarted one degree lower.
\end{itemize}\medskip
\ifdefined\hypertarget\hypertarget{sourcepointer-1}{}\fi
\noindent\textit{Source reference.} Original definition 014E: focused-seeds/AGP/sources/stacks-snapshot-20261011/tags.txt, source lines 1469--1469. Cross-chapter pointer into the Derived Categories chapter fixing the cone of a morphism of complexes; the induction step of the counted proof builds that cone itself and every cohomological conclusion drawn from it stays inside the counted proof span.\par\medskip
\ifdefined\hypertarget\hypertarget{sourcepointer-2}{}\fi
\noindent\textit{Source reference.} Original lemma 064R: sources/stacks-more-algebra/pinned-original-source.tex, source lines 15910--15910. Same-chapter pointer used once at proof line 16110 to transfer pseudo-coherence from the two outer terms of a distinguished triangle to the cone; it is one application of a standard two-out-of-three statement and carries no part of the induction that is not written out in the counted span.\par\medskip
\ifdefined\hypertarget\hypertarget{sourcepointer-3}{}\fi
\noindent\textit{Source reference.} Original lemma 064S: sources/stacks-more-algebra/pinned-original-source.tex, source lines 15965--15965. Same-chapter pointer used once at proof line 16111 to conclude that a single cohomology module of the cone is a finite R-module; the counted proof then exhibits the surjection onto that module explicitly.\par\medskip
\begin{lemma}
\label{lemma-pseudo-coherent}
Let $R$ be a ring. Let $K^\bullet$ be a complex of $R$-modules.
The following are equivalent
\begin{enumerate}
\item $K^\bullet$ is pseudo-coherent,
\item $K^\bullet$ is $m$-pseudo-coherent for every $m \in \mathbf{Z}$, and
\item $K^\bullet$ is quasi-isomorphic to a bounded above complex of finite
projective $R$-modules.
\end{enumerate}
If (1), (2), and (3) hold and $H^i(K^\bullet) = 0$ for $i > b$, then
we can find a quasi-isomorphism $F^\bullet \to K^\bullet$ with
$F^i$ finite free $R$-modules and $F^i = 0$ for $i > b$.
\end{lemma}

\begin{proof}
We see that (1) $\Rightarrow$ (3) as a finite free module is a finite
projective $R$-module. Conversely, suppose $P^\bullet$ is a bounded
above complex of finite projective $R$-modules. Say $P^i = 0$ for
$i > n_0$. We choose a direct sum decompositions
$F^{n_0} = P^{n_0} \oplus C^{n_0}$ with $F^{n_0}$ a finite free
$R$-module, and inductively
$$
F^{n - 1} = P^{n - 1} \oplus C^n \oplus C^{n - 1}
$$
for $n \leq n_0$ with $F^{n_0}$ a finite free $R$-module. As a complex
$F^\bullet$ has maps $F^{n - 1} \to F^n$ which agree with $P^{n - 1} \to P^n$,
induce the identity $C^n \to C^n$, and are zero on $C^{n - 1}$. The map
$F^\bullet \to P^\bullet$ is a quasi-isomorphism (even a homotopy equivalence)
and hence (3) implies (1).

\medskip\noindent
Assume (1). Let $E^\bullet$ be a bounded above complex of finite free
$R$-modules and let $E^\bullet \to K^\bullet$ be a
quasi-isomorphism. Then the induced maps
$\sigma_{\geq m}E^\bullet \to K^\bullet$ from the stupid truncation
of $E^\bullet$ to $K^\bullet$ show that $K^\bullet$ is $m$-pseudo-coherent.
Hence (1) implies (2).

\medskip\noindent
Assume (2). Since $K^\bullet$ is $0$-pseudo-coherent we see in particular
that $K^\bullet$ is bounded above. Let $b$ be an integer such that
$H^i(K^\bullet) = 0$ for $i > b$. By descending induction on
$n \in \mathbf{Z}$ we are going to construct finite free $R$-modules
$F^i$ for $i \geq n$, differentials $d^i : F^i \to F^{i + 1}$ for
$i \geq n$, maps $\alpha : F^i \to K^i$ compatible with differentials,
such that (1) $H^i(\alpha)$ is an isomorphism for $i > n$ and surjective for
$i = n$, and (2) $F^i = 0$ for $i > b$. Picture
$$
\xymatrix{
& F^n \ar[r] \ar[d]^\alpha & F^{n + 1} \ar[d]^\alpha \ar[r] & \ldots \\
K^{n - 1} \ar[r] & K^n \ar[r] & K^{n + 1} \ar[r] & \ldots
}
$$
The base case is $n = b + 1$ where we can take $F^i = 0$ for all $i$.
Induction step. Let $C^\bullet$ be the cone on $\alpha$
(Derived Categories, Definition \ref{derived-definition-cone}).
The long exact sequence of cohomology
$$
0 \to H^{n - 1}(K^\bullet) \to H^{n - 1}(C^\bullet) \to
H^n(F^\bullet) \to H^n(K^\bullet) \to H^n(C^\bullet) \to \ldots
$$
shows that $H^i(C^\bullet) = 0$ for $i \geq n$.
By Lemma \ref{lemma-cone-pseudo-coherent} we see that $C^\bullet$
is $(n - 1)$-pseudo-coherent. By Lemma \ref{lemma-finite-cohomology}
we see that $H^{n - 1}(C^\bullet)$ is a finite $R$-module.
In particular, we see that the kernel of $H^n(F^\bullet) \to H^n(K^\bullet)$
is a finite $R$-module. Choose a finite free $R$-module
$F^{n - 1}$ and a map $F^{n - 1} \to \Ker(F^n \to F^{n - 1})$ such that
$F^{n - 1}$ surjects onto the kernel of $H^n(F^\bullet) \to H^n(K^\bullet)$.
We extend our of complexes to
$$
\xymatrix{
& F^{n - 1} \ar[d]^{\alpha^{n - 1}} \ar[r] &
F^n \ar[r] \ar[d]^\alpha &
F^{n + 1} \ar[d]^\alpha \ar[r] & \ldots \\
\ldots \ar[r] &
K^{n - 1} \ar[r] & K^n \ar[r] & K^{n + 1} \ar[r] & \ldots
}
$$
Note that $\alpha^{n - 1}$ exists because the image of
$F^{n - 1} \to F^n \to K^n$ is in the image of $K^{n - 1} \to K^n$ by
construction. Dnote again $C^\bullet$ the cone of this extended
map of complexes. At this point we see that we get an exact sequence
$$
H^{n - 1}(F^\bullet) \to H^{n - 1}(K^\bullet) \to H^{n - 1}(C^\bullet) \to
H^n(F^\bullet) \cong H^n(K^\bullet) \to H^n(C^\bullet) \to \ldots
$$
In other words, we see that the cokernel of $H^{n - 1}(F^\bullet) \to
H^{n - 1}(K^\bullet)$ is a finite $R$-module, say generated by
the classes of $\xi_1, \ldots, \xi_r \in \Ker(K^{n - 1} \to K^n)$.
Then we replace $F^{n - 1}$ by $F^{n - 1} \oplus R^{\oplus r}$ where
the basis elements in the free summand map to zero in $F^n$ and the
to $\xi_i$ in $K^{n - 1}$. This
finishes the proof of the induction step.
\end{proof}

\end{document}
